\capitulo{2}{Objetivos del proyecto}

El objetivo principal del proyecto es diseñar e implementar un sistema de \eng{Retrieval-Augmented Generation} (\eng{RAG}) que permita enriquecer las respuestas de un modelo de lenguaje mediante la recuperación y utilización de información contextual proveniente de una colección documental. 

\section{Objetivos específicos}
Los objetivos identificados en el presente trabajo fin de grado son los siguientes:
\begin{itemize}
\tightlist
\item Recopilar la colección documental.
	\begin{itemize}
		\item Automatizar la obtención de los documentos (licitaciones del estado) a través de \eng{Web Scraping}.
	\end{itemize}
\item Preprocesar la colección documental.
	\begin{itemize} 
		\item Convertir la colección documenta (recopilada en \file{PDF}) a \file{Markdown}.
		\item Normalizar y limpiar el texto.
		\item Segmentar los documentos en fragmentos (\eng{chunks}).
		\item Generar \eng{embeddings} para representar semánticamente los fragmentos.
	\end{itemize}
	\item Implementar y gestionar un índice vectorial.
	\begin{itemize}
		\item Almacenar los \eng{embeddings} en una base de datos vectorial.
		\item Optimizar consultas de similitud para recuperar contexto relevante.
	\end{itemize}
\item Desarrollar la capa de interacción con el usuario
	\begin{itemize}
		\item Transformar la consulta en su representación vectorial.
		\item Recuperar los fragmentos más relevantes de la base de datos vectorial.
		\item Utilizar los fragmentos recuperados para enriquecer el \eng{prompt} del usuario y así construir el \eng{prompt} aumentado que finalmente se pasa al \eng{Large Language Model} (\eng{LLM}).
	\end{itemize}
\item Desarrollar interfaz que facilite al usuario la interacción con el modelo para que el proceso de aumentado del \eng{prompt} le sea transparente.
	\begin{itemize}
		\item Desarrollar una aplicación con \eng{Flask} para que el usuario pueda interactuar con el modelo.
		\item Levantar la aplicación \eng{Flask} con \eng{Docker} para asegurar el correcto funcionamiento de la aplicación en cualquier entorno.
		\item Desplegar la aplicación \eng{Flask} en un entorno preparado para la producción, aplicando \eng{Gunicorn} y \eng{Nginx}, garantizando la seguridad mediante un certificado SSL/TLS.
		\item Asignar un dominio DNS a la aplicación para facilitar el acceso. 
	\end{itemize}
\end{itemize}

\section{Objetivos personales}

\begin{itemize}
\tightlist
\item Adquirir conocimientos sobre los sistemas \eng{RAG}, comprendiendo sus fundamentos, componentes y aplicaciones en el ámbito de la inteligencia artificial.
\item Profundizar en el estudio de los modelos de lenguaje de gran tamaño (\eng{LLM}).
\item Desarrollar habilidades prácticas en la implementación de sistemas basados en \eng{RAG} y \eng{LLM}.
\item Aprender a aplicar los sistemas \eng{RAG} y los \eng{LLM} en la resolución de problemas reales.
\item Aplicar los conocimientos adquiridos durante la carrera.
\end{itemize}
